% source: https://github.com/pfradin/luadraw/discussions/140#discussioncomment-15060971 (pfradin, block 1)
\documentclass[boder=5pt]{standalone}
\usepackage[svgnames]{xcolor}
\usepackage[3d]{luadraw}
\usepackage{fourier-otf}
\begin{document}

\colorlet{myOrange}{orange!35}%

\begin{luadraw}{name=tubes}
local g = graph3d:new{window={-6,6,-4,7},bg="lightgray!30" }
Hiddenlines = false; Hiddenlinestyle = "noline"
g:Linewidth(2); g:Linecolor("Orange")
local O, r = Origin, 0.5
local R  = r*math.sqrt(2)
local S = sphere(O,R)
local C = parallelep(M(-r,-r,-r), 2*r*vecI,2*r*vecJ,2*r*vecK)
local crossing_base = clip3d(S,C)

local crossing = function(center)
    return shift3d(crossing_base,center)
end

local Shift = function(L,u)
    local f = function(A)
        if isPoint3d(A) then return A+u else return A end
    end
    return ftransform3d(L,f)
end

local grid_base = {} -- list of {center, "type"}
for y = -4, 4  do
    local A = y*2*r*vecJ
    if y%3 == 0 then
        table.insert(grid_base, {A,"c"})
        table.insert(grid_base, {A+2*r*vecI,"x"})
        table.insert(grid_base, {A-2*r*vecI,"x"})
        table.insert(grid_base, {A+2*r*vecK,"z"})
        table.insert(grid_base, {A-2*r*vecK,"z"})
    else
        table.insert(grid_base, {A,"y"})
    end
end

local k = 6*r
local grid = concat(grid_base,Shift(grid_base,k*vecI),Shift(grid_base,-k*vecI))
grid = concat(grid, Shift(grid,k*vecK))
table.sort(grid, function(e1,e2) return pt3d.dot(e1[1],g.Normal) < pt3d.dot(e2[1],g.Normal) end)

for _,elt in ipairs(grid) do
    local A, mode = table.unpack(elt)
    if mode == "c" then -- crossing
        g:Dfacet(crossing(A), {color="myOrange",contrast=0,backcull=true,edgecolor="myOrange",mode=mShaded})
    elseif mode == "x" then
        g:Dcylinder(A-r*vecI,r,A+r*vecI, {color="orange"})
    elseif mode == "y" then
        g:Dcylinder(A-r*vecJ,r,A+r*vecJ, {color="orange"})
    elseif mode == "z" then
        g:Dcylinder(A-r*vecK,r,A+r*vecK, {color="orange"})
    end
end 

g:Show()
\end{luadraw}
\end{document} 
